\documentclass[12pt]{article}
\usepackage[a3paper, margin=1in]{geometry}
\usepackage{amsmath, amssymb, amsthm, ulem, booktabs}

\begin{document}

\section*{Cvičení 7. 5. 2026}

Vypracoval: Daniel Ransdorf \hfill Podpis: \rule{4cm}{0.4pt}


\bigskip

\section{Problém}
Dva hráči na ploše 3x3 střídavě zapisují své značky - O nebo X (písmeno "O" nebo písmeno "X"). Pokud se některému hráči podaří umísit tři své značky v řadě - vodorovně, svisle nebo úhlopříčně, vyhrál. Známe počáteční stav, chceme vyhodnotit, zda hráč na tahu může vynutit výhru, remízu, nebo jestli nuceně prohrává.
\section{Návrh}
Po celý návrh ignorujme neplatné stavy, kde mají výherní trojici oba hráči (v programování toto vyrazíme validací).
Pozorujme:
\begin{enumerate}
  \item mezi stavama piškvorek neexistují cykly (každý nový stav má o volné pole méně)
  \item každý stav vede do konečného stavu maximálně v 9 krocích (ze stejného důvodu)
\end{enumerate}

\subsection{Základní definice}

Označme aktuální stav $S$, BÚNO zvolme X ("křížky") jako hráče na tahu. Označme $S_1,S_2,...,S_n$ herní stavy, po zahrání tahu ve stavu $S$.
Označme funkci $f_X$, která vyhodnocuje, zda je stav z pohledu X nucená výhra/remíza/prohra. Definujme hodnoty výhra=1,remíza=0,prohra=-1,
s ideou, abychom měli uspořádání stavů (výhra>remíza>prohra).
Předpokládáme, že X hraje nejlepší možné tahy, čili výsledkem aktuálního stavu je výsledek nejlepšího tahu. Definujme:

\[
  f_X(S) := max\left\{ f_X(S_1),...,f_X(S_n) \right\}
\]

Aby byla definice kompletní, mějme stav $A$, kde je na tahu hráč O, označme následné stavy $A_1,...,A_n$.
Hráč O hraje taky nejlepší tahy, čili z našeho pohledu je výsledek "nejhorší možný, co O může zahrát".
Definujme tedy:

\[
  f_X(A) := min\left\{ f_X(A_1),...,f_X(A_n) \right\}
\]

Ještě triviálně zadefinujeme všechny konečné stavy, aby rekurze někde skončila. Z pozorování víme,
že rekurze někdy skončí, také, že se nemá, jak zacyklit.

(obdobně zavedeme $f_O$)

\subsection{Zjednodušení pro implementaci}

Nechceme rozlišovat, ze kterého pohledu výsledky hodnotíme. Zaveďme jen jednu funkci $eval$, která stav hodnotí
z pohledu hráče na tahu. Když bude na tahu druhý hráč (ne ten z našeho pohledu), stačí otočit znaménko.
Definice funkce pak bude vypadat podobnou úvahou takto:

\[
  f(S) := max\left\{ -f(S_1),...,-f(S_n) \right\}
\]

Zjednodušme na:

\[
  f(S) := -min\left\{ f(S_1),...,f(S_n) \right\}
\]

\subsection{Technikálie}

Výsledek stavů bychom mohli memoizovat. Jakmile nám v rekurzi nějaká funkce zvýší výsledek na $1$, můžeme ihned vracet,
není nic lepšího, co hledat.

\end{document}
